% Generated by tex/build_swebok_structure.py from tex/chapters/ch01/09_実務上の考慮.tex
\subsubsection{要求プロセスの反復性}

現実の要求活動は、一回で終わらない。最初は広く浅く要求を集め、次に重要な部分を深く分析し、その後また不足が見つかって追加で獲得することが多い。要求は幅と深さの両方を持つため、反復的に進めるのが自然である。

ウォーターフォール型では、要求フェーズで多くの要求を先に固める。反復型やアジャイル型では、初期に高レベル要求を押さえ、その後、開発する順に詳細化する。重要なのは、ライフサイクルの名前ではなく、対象ソフトウェアの性質に合う要求活動を選ぶことである。
